Route choice is not only a shortest-path problem. A person may prefer a greener, quieter, simpler, better-lit, or more familiar route even when that route is not the minimum-distance option. This makes routing a natural personalization problem: the best route depends on the origin and destination, but also on the request context and the traveler's implicit or explicit preferences. Prior work on personalized route recommendation often assumes large-scale navigation logs, clean historical route choices, or trained models over dense trajectory datasets~\cite{dai2015bigtrajectory, nasr2019, userhabits2024}. This paper targets a complementary open-data setting: interpretable route-candidate reranking that can be built, inspected, and reproduced with OSM features and public GPS-derived movement signals.

This paper studies the following question: can OSM-derived historical movement signals support preference-aware route-candidate reranking when clean user histories are unavailable? We use the term \emph{pseudo-history profile} because public GPS traces provide movement evidence, not guaranteed persistent user identity. That identity gap is the core research challenge. The system therefore treats trace reconstruction quality as a measured component of the pipeline and evaluates route-candidate reranking separately from trace reconstruction.

We present a working OSM-based route-candidate reranking system. Given an origin, destination, request context, optional preference text, and optional history records, the backend generates OSM route candidates, extracts interpretable route features, constructs dynamic profile weights from contextual history, and reranks candidates using profile, prompt, or hybrid modes. The implementation is exposed through a FastAPI backend and React interface, preserving the product-facing demonstration that motivated the project. This matters because route selection is ultimately an interactive decision aid: the system must produce ranked alternatives and explanations that a user can inspect, not only offline scores.

The evaluation is designed around this validity challenge. A first OSM public-trace experiment demonstrates the end-to-end pipeline for Toronto walking traces and uses route-distance ratio to expose reconstruction cases that require improvement. This Toronto probe is exploratory internal validation, not the headline public benchmark. The primary public benchmark is Porto: it reconstructs observed OSM driving routes from the public Porto taxi trajectory dataset, generates route candidates for the same origin and destination, and reports ranking metrics plus path-overlap metrics with bootstrap confidence intervals. Walking OSM traces and Porto taxi trajectories represent different travel modes, so they are not merged into one headline result. The Porto setup creates a concrete bridge from an open-data reranking system to future same-data comparisons against learned route-search methods such as NASR~\cite{nasr2019}.

The contributions are:

\begin{itemize}
  \item An end-to-end OSM route-candidate reranking architecture that combines candidate generation, interpretable route features, contextual pseudo-history profiles, prompt-based reranking, and a GUI/API workflow.
  \item A trace-to-profile pipeline that converts public GPS-derived signals into pseudo-history records while reporting reconstruction quality instead of assuming clean user histories.
  \item A route-candidate evaluation protocol in which generated candidates are ranked against held-out observed routes, avoiding history-record ranking artifacts that can overstate shortest-distance baselines.
  \item A Porto public-dataset benchmark with random, shortest-distance, profile, trajectory-derived vehicle profile, learned feature-ranker, feature-oracle, and path-oracle comparisons using Hit@K, MRR, NDCG@3, feature distance, path F1, NDTW, edit distance, paired bootstrap differences, and confidence intervals.
\end{itemize}

The central claim is that profile-aware OSM route-candidate reranking can be evaluated transparently against non-personalized baselines in an open-data setting, and that the current Porto benchmark shows measurable ranking gains for trajectory-derived profiles while path metrics expose remaining limitations. The remaining research frontier is to strengthen map matching, candidate generation, and reproduced external comparisons so that the same framework can support a direct benchmark claim.
